\section{Variations of subgraph densities and relative entropy}\label{sec:density-variations}

We first express the first-order changes in subgraph densities through
their derivative kernels.  Bounds on these kernels allow us to
construct feasible perturbations and prove that every upper-tail minimizer
satisfies $p\le W\le1-\gamma$ for some $\gamma>0$.  With this interiority,
we can then differentiate the relative entropy in any bounded symmetric direction.

\subsection{Derivative kernels of subgraph densities}\label{sec:derivatives}

Let $W$ be a graphon and $V\in L^\infty_{\mathrm{sym}}([0,1]^2)$.
The integral formulas in~\eqref{eq:graphon-quantities} extend $e$ and
$t(H,\cdot)$ to $L^\infty_{\mathrm{sym}}([0,1]^2)$ as linear and
polynomial functionals, respectively. In particular, $t(H,W+\zeta V)$
is polynomial in $\zeta\in\RR$, even when $W+\zeta V$ is not a graphon.
We therefore define
the \emph{directional derivative} of $t(H,\cdot)$
at $W$ in the direction $V$ by
\[
  Dt(H,W)[V]:=
  \left.\frac{\dd}{\dd\zeta}t(H,W+\zeta V)\right|_{\zeta=0}.
\]
This derivative describes the first-order change in the $H$-density under the
perturbation $W+\zeta V$.  In the integral defining $t(H,W+\zeta V)$,
each edge $ij$ contributes the factor
$W(x_i,x_j)+\zeta V(x_i,x_j)$.  The product rule therefore gives
\[
  Dt(H,W)[V]
  =\sum_{ab\in E(H)}
    \int_{[0,1]^{V(H)}}V(x_a,x_b)
    \prod_{ij\in E(H)\setminus\{ab\}}W(x_i,x_j)
    \prod_{i\in V(H)}\dd x_i.
\]

For each edge $ab$ of $H$, choose an ordering $(a,b)$ of its endpoints
and fix $x_a=x$ and $x_b=y$, where $x,y\in[0,1]$.  Integrating the remaining
edge factors over all other vertex labels defines
\[
  K_{ab,W}(x,y):=\int_{[0,1]^{V(H)\setminus\{a,b\}}}
  \prod_{ij\in E(H)\setminus\{ab\}}W(x_i,x_j)
  \prod_{i\in V(H)\setminus\{a,b\}}\dd x_i,
  \qquad x_a=x,\ x_b=y.
\]
The derivative then takes the form
\[
  Dt(H,W)[V]=\sum_{ab\in E(H)}
  \int_{[0,1]^2}K_{ab,W}(x,y)V(x,y)\dd x\dd y.
\]
Reversing the endpoint order gives $K_{ba,W}(x,y)=K_{ab,W}(y,x)$.
Since $V$ is symmetric, exchanging the integration variables shows that
the contribution from each edge is independent of its orientation.
Averaging over the two orientations gives the symmetric kernel
\begin{equation}\label{eq:derivative-kernel}
  \mathcal D_{H,W}(x,y):=\frac12\sum_{ab\in E(H)}\bigl(K_{ab,W}(x,y)+K_{ab,W}(y,x)\bigr),
\end{equation}
called the \emph{derivative kernel} of the $H$-density at $W$.
For every $V\in L^\infty_{\mathrm{sym}}([0,1]^2)$, it satisfies
\begin{equation}\label{eq:density-derivative}
  Dt(H,W)[V]=\int_{[0,1]^2}\mathcal D_{H,W}(x,y)\,V(x,y)\dd x\dd y.
\end{equation}
Since each $K_{ab,W}$ takes values in $[0,1]$, we have
$0\le\mathcal D_{H,W}\le m$.
If, in addition, $W\ge\omega>0$ almost everywhere, then
$\mathcal D_{H,W}\ge m\omega^{m-1}$.

For graphons $V,W$, apply~\eqref{eq:product-difference} to the $m-1$
edge factors defining $K_{ab,V}$ and $K_{ab,W}$. All other factors lie
in $[0,1]$, so integrating over all vertex labels bounds each summand's
contribution to the $L^1([0,1]^2)$ norm by $\|V-W\|_1$. Thus
\[
  \|K_{ab,V}-K_{ab,W}\|_1\le(m-1)\|V-W\|_1.
\]
Summing over the $m$ edges and averaging over their two orientations
in~\eqref{eq:derivative-kernel} gives
\begin{equation}\label{eq:derivative-kernel-lipschitz}
  \|\mathcal D_{H,V}-\mathcal D_{H,W}\|_1
  \le m(m-1)\|V-W\|_1.
\end{equation}

\begin{expl}\label{expl:triangle-derivative}
Let $H=K_3$ have vertices $1,2,3$.  For the edge $12$, fixing
$x_1=x$ and $x_2=y$ leaves the factors $W(x,x_3)W(y,x_3)$.
The other two edges give the same kernel, so
\[
  K_{12,W}(x,y)=K_{13,W}(x,y)=K_{23,W}(x,y)
  =\int_0^1W(x,z)W(y,z)\dd z.
\]
This kernel is symmetric in $x,y$.  Summing the three edge contributions
in~\eqref{eq:derivative-kernel} therefore gives
\[
  \mathcal D_{K_3,W}(x,y)=3\int_0^1W(x,z)W(y,z)\dd z.
\]
\end{expl}

\begin{expl}\label{expl:path-derivative}
Let $H=P_3$ have vertices $1,2,3$ and edges $12,23$.
Choosing the endpoint orders $(1,2)$ and $(2,3)$ gives
\[
  K_{12,W}(x,y)=\int_0^1W(y,z)\dd z=d_W(y),\qquad
  K_{23,W}(x,y)=\int_0^1W(z,x)\dd z=d_W(x).
\]
For each edge, averaging over its two orientations gives
$\tfrac12(d_W(x)+d_W(y))$.  The two edges therefore contribute
\[
  \mathcal D_{P_3,W}(x,y)=d_W(x)+d_W(y).
\]
\end{expl}

\subsection{Interiority and entropy derivatives}\label{sec:upper-tail-interiority}

Fix $p$, $H_1,\dots,H_s$, and $r_1,\dots,r_s$ as in
\Cref{thm:joint-upper-tail}, write $t_j(U):=t(H_j,U)$ for every
$U\in L^\infty_{\mathrm{sym}}([0,1]^2)$,
and let $W$ minimize~\eqref{eq:joint-variational-problem}.
An index $j$ is \emph{active} if $t_j(W)=r_j^{m_j}$,
and denote by $\mathcal A$ the set of active indices.

\begin{lemma}\label{lem:upper-tail-basic}
Let $W$ be a minimizer of~\eqref{eq:joint-variational-problem}.
If the constant graphon $p$ is not feasible, the active set $\mathcal A$
is nonempty.
In all cases, there is $\gamma>0$ such that $p\le W\le1-\gamma$ almost everywhere.
\end{lemma}

\begin{proof}
If the constant graphon $p$ is feasible, it is the unique minimizer,
and the claimed bounds hold with $\gamma=1-p$.  We may therefore assume
that $p$ is not feasible, so $W\not\equiv p$ and $I_p(W)>0$.

We first show that $\mathcal A$ is nonempty. If it were empty, then
$t_j(W)>r_j^{m_j}$ for every $j$.
Put $W_\lambda:=p+\lambda(W-p)$.  Each $\lambda\mapsto t_j(W_\lambda)$ is continuous,
so we can choose $\lambda<1$ arbitrarily close to $1$ such that $t_j(W_\lambda)>r_j^{m_j}$ for every $j$.
Convexity of $J_p$ and $J_p(p)=0$ then give a contradiction to the minimality of $W$:
\[
  I_p(W_\lambda)\le \lambda I_p(W)+(1-\lambda)J_p(p)= \lambda I_p(W)<I_p(W).
\]
Thus at least one constraint is active.

We next prove the lower bound $W\ge p$. Put $W^+:=\max(W,p)$.
Then $t_j(W^+)\ge t_j(W)$ for every $j$, so $W^+$ is
feasible.  Since $J_p$ is strictly decreasing on $[0,p]$, we have
$I_p(W^+)\le I_p(W)$, with equality if and only if $W\ge p$ almost everywhere.
Since $W$ is a minimizer, it follows that $W\ge p$.

It remains to show that $W$ is essentially bounded away from $1$;
we first outline the competitor construction for this upper bound.
The graphon $W$ cannot equal $1$ almost everywhere: otherwise every density
would equal $1>r_j^{m_j}$, contradicting the existence of an active
constraint.  We can therefore fix $\gamma_0>0$ and a symmetric set $A$
of positive measure on which $p\le W\le1-\gamma_0$.  Suppose, for a
contradiction, that $W$ is not essentially bounded away from $1$.
For any sufficiently large $\Lambda$ with
$e^{-\Lambda}<\min\{\gamma_0,1-p\}$, the symmetric set
\[
  E:=\{W>1-e^{-\Lambda}\}
\]
has positive measure and is disjoint from $A$. We will lower $W$ on $E$
and compensate by raising it on $A$ so that every density constraint
remains satisfied. Lowering $W$ on $E$ gives a large entropy saving per
unit decrease, since $J_p'$ diverges at $1$. On $A$, the entropy cost per
unit increase stays bounded. We will choose the perturbation so that
the saving on $E$ outweighs the cost on $A$.

We now choose the compensating increase on $A$ and verify that the
resulting perturbation is feasible.
Since $W\ge p>0$ and both $E$ and $A$ have positive measure, the
integrals of each derivative kernel over these sets are positive.
Choose
\[
  c:=2\max_{1\le j\le s}
       \frac{\int_E\mathcal D_{H_j,W}}{\int_A\mathcal D_{H_j,W}}>0,
  \qquad
  Z:=-\mathbf1_E+c\mathbf1_A.
\]
The derivative formula~\eqref{eq:density-derivative} for the $H_j$-density gives
\[
    Dt_j(W)[Z]
    =-\int_E\mathcal D_{H_j,W}+c\int_A\mathcal D_{H_j,W}
    \ge\int_E\mathcal D_{H_j,W}>0
    \qquad(1\le j\le s).
\]
Using this direction, define the competitor for small $\theta>0$ by
\[
  V_\theta:=W+\theta Z
  =
  \begin{cases}
    W-\theta, & \text{on }E,\\
    W+\theta c, & \text{on }A,\\
    W, & \text{outside }A\cup E.
  \end{cases}
\]
For fixed $\Lambda$, each density is polynomial in $\theta$, and its
strictly positive first derivative gives
\[
  t_j(V_\theta)
  =t_j(W)+\theta\,Dt_j(W)[Z]+O_\Lambda(\theta^2)>t_j(W)
\]
for all sufficiently small $\theta>0$.  Since there are only finitely
many constraints, each of them holds for the same sufficiently small
$\theta$.

We next check that $V_\theta$ takes values in $[0,1]$.  Shrink $\theta$
further so that
\[
  0<\theta\le1-e^{-\Lambda}-p,
  \qquad
  \theta c\le\gamma_0/2.
\]
These conditions are possible because $\Lambda$ and $c$ are fixed
and $1-e^{-\Lambda}>p$.  They give the following bounds
on the indicated sets:
\[
  p\le1-e^{-\Lambda}-\theta<W-\theta\le1\quad\text{on }E,\qquad
  p\le W+\theta c\le1-\gamma_0/2\quad\text{on }A.
\]
Outside $A\cup E$, the values are unchanged.  Since the sets and the
perturbation are symmetric, $V_\theta$ is a graphon and hence a
feasible competitor.

With feasibility established, it remains to prove $I_p(V_\theta)<I_p(W)$
by comparing the contributions from $E$ and $A$. On $E$, we have
$V_\theta=W-\theta$. We may shrink $\theta$ further to ensure that
$0<\theta\le e^{-\Lambda}/2$.  For $1-e^{-\Lambda}\le z\le1$ and $q\in[z-\theta,z)$, we have
\[
  q\ge z-\theta\ge1-e^{-\Lambda}-\theta\ge p,
  \qquad
  1-q\le1-z+\theta\le\tfrac32e^{-\Lambda}.
\]
Consequently, \eqref{eq:entropy-derivatives} gives
\[
    J_p'(q)
    =\log\frac qp+\log\frac{1-p}{1-q}
    \ge\log\frac{2(1-p)}{3e^{-\Lambda}}
    =\Lambda-C_p,
    \qquad C_p:=\log\frac{3}{2(1-p)}.
\]
Integrating over $[z-\theta,z]$ gives
\[
  J_p(z-\theta)-J_p(z)
  =-\int_{z-\theta}^zJ_p'(q)\dd q
  \le-\theta(\Lambda-C_p).
\]
The endpoint case $z=1$ follows by continuity of $J_p$.  Applying this
bound pointwise on $E$ yields
\[
  \int_E\bigl(J_p(V_\theta)-J_p(W)\bigr)
  \le-\theta(\Lambda-C_p)\nu(E).
\]
Thus, relative to $\theta\nu(E)$, the cost saving on $E$ grows without bound
as $\Lambda$ increases.

On $A$, both $W$ and $V_\theta$ take values in the fixed interval
$[p,1-\gamma_0/2]$.  Let $C$ bound $|J_p'|$ on this interval;
$C$ depends only on $p$ and $\gamma_0$.
To control the total cost on $A$, we bound $c\nu(A)$ in terms of $\nu(E)$.
Since $W\ge p$, we have
\[
  m_jp^{m_j-1}\le\mathcal D_{H_j,W}\le m_j.
\]
It follows that
\[
  \frac{\int_E\mathcal D_{H_j,W}}{\int_A\mathcal D_{H_j,W}}
  \le\frac{m_j\nu(E)}{m_jp^{m_j-1}\nu(A)}
  =p^{1-m_j}\frac{\nu(E)}{\nu(A)}.
\]
By taking the maximum over $j$, we obtain
\[
  c\nu(A)\le2p^{1-\max_jm_j}\nu(E).
\]
Integrating along the increase on $A$ and using this bound gives
\[
    \int_A\bigl(J_p(V_\theta)-J_p(W)\bigr)
    =\int_A\int_W^{W+\theta c}J_p'(q)\dd q
    \le C\theta c\nu(A)
    \le2C\theta p^{1-\max_jm_j}\nu(E),
\]
so the cost increase on $A$ is bounded by a constant multiple of
$\theta\nu(E)$, independently of $\Lambda$.

Combining the two contributions gives
\begin{align*}
  I_p(V_\theta)-I_p(W)
  &=\int_E\bigl(J_p(V_\theta)-J_p(W)\bigr)
    +\int_A\bigl(J_p(V_\theta)-J_p(W)\bigr)\\
  &\le-\theta\bigl(\Lambda-C_p-2Cp^{1-\max_jm_j}\bigr)\nu(E).
\end{align*}
The constants $C_p$ and $C$ are independent of $\Lambda$.  Choose
$\Lambda$ large enough to make the coefficient in parentheses
positive, then take $\theta>0$ sufficiently small to satisfy all
the preceding feasibility conditions.  Since $\nu(E)>0$, the resulting
feasible graphon $V_\theta$ has strictly smaller cost, contradicting
the minimality of $W$.
\end{proof}

By \Cref{lem:upper-tail-basic}, $W$ takes values in a compact subinterval
of $(0,1)$.  For every bounded symmetric direction $V$, the function
$W+\zeta V$ is therefore a graphon for sufficiently small $|\zeta|$,
and differentiating the entropy under the integral gives
\begin{align*}
  DI_p(W)[V]
  &:=\left.\frac{\dd}{\dd\zeta}I_p(W+\zeta V)\right|_{\zeta=0} \\
  &=\int_{[0,1]^2}\left.\frac{\dd}{\dd\zeta}J_p\bigl(W(x,y)+\zeta V(x,y)\bigr)\right|_{\zeta=0}\dd x\dd y \\
  &=\int_{[0,1]^2}J_p'(W(x,y))V(x,y)\dd x\dd y.
\end{align*}
